\subsection*{语法速览}

\begin{table}[htbp]\centering\small
\begin{tabularx}{\linewidth}{@{}lX@{}}
\toprule
\textbf{语法} & \textbf{含义} \\
\midrule
 &  \\
\texttt{dma.copy.swiz\textless{}N\textgreater{}} src =\textgreater{} dst & 采用以字节计的重排宽度 \texttt{N} 的拷贝 \\
\texttt{mma.load.swiz\textless{}N\textgreater{}} src & 与重排 \texttt{N} 一致的 MMA 操作数加载 \\
\texttt{.view(M, N).from(r, c)} & 任意偏移处的 \texttt{M x N} 窗口 \\
\texttt{.subspan(M, K).step(sM, sK).at(i, j)} & 带步长的分块选择 \\
\texttt{.zfill} & 把拷贝源中的越界元素补零 \\
\texttt{span\_as(dims)} & 把线性存储重新解释为带形状的张量 \\
\bottomrule
\end{tabularx}
\end{table}

3.6 节把上面的分块表达式用作搬运语句的源与目的。对于程序中确实需要纯 C++ 的部分，2.13 节汇集了 \co@__device__@ 函数、预处理守卫与注入的源代码文本。
